\documentclass[11pt,a4paper]{article}
% Shared preamble of the English and German IT security guides (guide.tex, guide-de.tex).
\usepackage[utf8]{inputenc}
\usepackage[T1]{fontenc}
\usepackage{lmodern}
\usepackage{microtype}
\usepackage[margin=0.9in]{geometry}
\usepackage{booktabs}
\usepackage{tabularx}
\usepackage{array}
\usepackage{enumitem}
\usepackage{fontawesome5}
\usepackage{xcolor}
\usepackage{listings}
\usepackage{tcolorbox}
\usepackage[colorlinks=true,linkcolor=blue,citecolor=blue,urlcolor=blue]{hyperref}
\setlength{\emergencystretch}{2em}
\setlist{itemsep=2pt,topsep=3pt}

\definecolor{separblue}{HTML}{0550AE}
\definecolor{separbg}{HTML}{F6F8FA}
\definecolor{separborder}{HTML}{D0D7DE}
\definecolor{separwarn}{HTML}{FFF8C5}
\definecolor{separwarnborder}{HTML}{D4A72C}

\lstset{
  backgroundcolor=\color{separbg}, frame=single, rulecolor=\color{separborder},
  basicstyle=\ttfamily\small, breaklines=true, columns=fullflexible, keepspaces=true,
  showstringspaces=false, aboveskip=6pt, belowskip=6pt, xleftmargin=4pt, xrightmargin=4pt,
  upquote=true
}
\newtcolorbox{callout}[1][]{colback=separbg, colframe=separblue, boxrule=0.8pt, arc=3pt,
  left=8pt, right=8pt, top=6pt, bottom=6pt, #1}
\newtcolorbox{warnbox}[1][]{colback=separwarn, colframe=separwarnborder, boxrule=0.8pt, arc=3pt,
  left=8pt, right=8pt, top=6pt, bottom=6pt, #1}

% a fact the repository does not hold yet (listed in the build report)
\newcommand{\todo}[1]{\textcolor{red}{\textbf{[TODO:} #1\textbf{]}}}
\newcommand{\cmd}[1]{\texttt{#1}}
\newcommand{\contact}{\href{mailto:parsa@parsabe.com}{\texttt{parsa@parsabe.com}}}
\newcommand{\repo}{\href{https://github.com/parsabe/Separ}{\texttt{github.com/parsabe/Separ}}}
\newcommand{\releases}{\href{https://github.com/parsabe/Separ/releases}{\texttt{github.com/parsabe/Separ/releases}}}


\title{\Large\bfseries Separ v2.4 on TU Bergakademie Freiberg Linux machines\\[4pt]
\large Guide for the IT security officer, URZ administrators and the data protection officer}
\author{Parsa Besharat (student research project)\\ \contact \quad|\quad \repo}
\date{November 2026 \quad|\quad applies to Separ 2.4.0}

\begin{document}
\maketitle

\begin{callout}
\textbf{How to read this guide.} Section~\ref{sec:summary} is a one-page summary for
decision-makers. Sections~\ref{sec:does}--\ref{sec:approvals} explain what Separ does and which
approvals are needed. Sections~\ref{sec:req}--\ref{sec:pilot} are for the administrators who run
the pilot. Every command and option in this guide is taken from \texttt{docs/it-facts.md} in the
repository, which records the exact help text of Separ 2.4.0. Network ranges, names and addresses
in examples are placeholders such as \texttt{<your campus range>}.
\end{callout}

\tableofcontents
\newpage

% ============================================================================
\section{Summary for decision-makers}
\label{sec:summary}

\paragraph{What Separ is.} A student research project (version 2.4.0) that looks at actions on a
Linux machine and asks three questions: did the input come from outside (a web page, a mail, a
download, an AI agent's input, a request to a web service)? Is the program doing something its
role never does (a web server starting a shell)? Does it try something dangerous (download and run
code, hide, start itself at boot, read passwords)? For commands that are routed through it, Separ
decides \emph{before} they run: allow, hold (ask), or block. For sensor data an administrator
connects, it decides after the fact. Separ contains no machine learning and sends nothing to the
internet.

\paragraph{What is asked.} An \textbf{observe-only pilot on 3 to 5 Linux machines for about two
weeks}. In observe mode Separ blocks nothing: it only records what it \emph{would} hold or block,
so that false alarms on real work can be counted before anyone considers enforcement.

\paragraph{What it records.} On each machine: one line per observed event and verdict, in which
commands, web addresses, file paths and file contents are replaced by keyed fingerprints
(HMAC-SHA256 with a random per-machine key). It does store some identifiers in clear text: login
names, host names reported by sensors, and names of quarantined files (Section~\ref{sec:facts}).
Nothing leaves the machine unless an administrator configures a notification channel; those
messages name the machine by pseudonym and contain no user name.

\paragraph{What it cannot do.}
\begin{itemize}
  \item It sees only what is routed to it: commands sent through \cmd{separ guard}, the browser
        extension, a watched folder, and connected sensors. A command pasted into a terminal
        without the guard, or a script run by name, is not seen before it runs.
  \item In block mode it also blocks legitimate install instructions pasted from a web page (in
        tests, 5 of 10 everyday pasted commands were blocked). That is why the pilot is observe-only.
  \item Containment of real processes happens only when the system service runs as root with
        \cmd{--enforce}; otherwise containment is simulated. Root, and on per-user setups the same
        user, can stop Separ.
  \item It has not been measured on real users or servers. CPU and memory use are unknown and are
        to be measured in the pilot.
\end{itemize}

\paragraph{Decision needed.} Approval by the IT security officer, a data protection decision by the
data protection officer, and, for staff machines or servers used by staff, the staff council
(Section~\ref{sec:approvals}).

% ============================================================================
\section{What Separ does and does not do}
\label{sec:does}

\subsection{The three questions}
\begin{description}
  \item[Origin.] Did the triggering input come from outside the trust boundary? Content from a web
    page, mail, download, removable medium, a request from the internet or an AI agent's input gets
    a \emph{label} that fades over 30 minutes unless the story turns dangerous.
  \item[Contract.] Is the program acting outside its role? Each role (interactive shell, web server
    worker, package installer, \ldots) has a signed contract of what it may do.
  \item[Effect.] Does the story try something dangerous: fetch from an unknown place, hide, persist,
    escalate privileges, move laterally, read credentials, tamper with security tools, change many
    files at once?
\end{description}
Separ adds up the evidence of the three questions into a score. Above one threshold it holds (asks),
above a higher one it blocks; it never blocks on suspicion alone without a contract violation.
Some events always block: reading a decoy (canary) password, a labelled program touching a
protected asset, and attempts to stop or alter Separ.

\subsection{How a decision is made, step by step}
\label{sec:flow}
Figure~\ref{fig:flow} follows one event through Separ. The numbers in the figure are the steps:
(1)~an event arrives from one of the places Separ watches; (2)~the command text is replaced by a
fingerprint, so nothing after the dashed red line ever sees it; (3)~content from outside is
labelled; (4)~related events become one story; (5)~the story is judged on the three questions;
(6)~an override decides at once, otherwise the score decides; (7)~the mode decides whether
anything happens; (8)~in block mode the story's own processes and files are contained; (9)~every
decision is logged and can be notified. In the pilot (observe mode) the path ends at step~7 with
``only recorded''.

\begin{figure}[!htbp]
  \centering
  \adjustbox{max width=\textwidth}{% Separ v2.4 end to end: where events come from, the stages, the decision, what the mode does,
% containment and reporting. One picture for the paper (main.tex) and both IT guides
% (it-service-desk/guide*.tex). Needs \input{diagrams/styles} (tikz, xcolor). Facts from the code:
% pipeline/pipeline.py (process), pipeline/judge.py, calculus/evidence.py (thresholds of the default
% policy), gate.py (modes, hold flows), contain.py, enforcer_linux.py (serve --enforce).
\begin{tikzpicture}[rm, x=1cm, y=1cm]
  % ---- 1: where Separ looks ------------------------------------------------------------------
  \foreach \x/\n/\t/\s in {
      1.7/in1/{\texttt{separ guard -{}- <cmd>}}/{before the command runs},
      5.15/in2/{Browser extension}/{clipboard writes, downloads},
      8.6/in3/{Watched folders}/{sandbox, web root},
      12.05/in4/{Connected sensors}/{Tetragon, Falco, sshd, IdP},
      15.5/in5/{Rootkit scan}/{hourly, read-only}} {
    \node[rm io, minimum width=30mm, minimum height=9mm] (\n) at (\x,0) {\t\\\tiny \s};
  }
  \node[rm group label] at (-0.05,0.62) {1\enspace WHERE SEPAR LOOKS};
  \node[rm raw, minimum width=172mm, minimum height=5mm] (raw) at (8.6,-1.3)
    {raw event, in memory only (at most 10 minutes, never written to disk)};
  \foreach \n in {in1,in2,in3,in4,in5} { \draw[rm flow] (\n.south) -- (\n.south |- raw.north); }

  % ---- 2-5: the stages ----------------------------------------------------------------------
  \node[rm stage, minimum width=37mm, minimum height=15mm] (obs) at (2.1,-3.15)
    {\textbf{2\enspace Observe}\\\tiny command, URL, path $\to$ keyed HMAC\\\tiny trusted host? canary? obfuscated?\\\tiny raw text dropped};
  \node[rm stage, minimum width=37mm, minimum height=15mm] (tag) at (6.4,-3.15)
    {\textbf{3\enspace Tag}\\\tiny outside content gets a label\\\tiny it spreads to child processes\\\tiny and files; half-life 30 min};
  \node[rm stage, minimum width=37mm, minimum height=15mm] (lnk) at (10.7,-3.15)
    {\textbf{4\enspace Link}\\\tiny related events $\to$ one story\\\tiny paste $\to$ curl $\to$ sh $\to$ file};
  \node[rm stage, minimum width=37mm, minimum height=15mm] (jdg) at (15.0,-3.15)
    {\textbf{5\enspace Judge: three questions}\\\tiny Origin: from outside?\\\tiny Contract: outside its role?\\\tiny Effect: dangerous action?};
  \draw[rm flow] (raw.south -| obs.north) -- (obs.north);
  \draw[rm flow] (obs) -- (tag);
  \draw[rm flow] (tag) -- (lnk);
  \draw[rm flow] (lnk) -- (jdg);
  \draw[rm wall] (4.25,-2.05) -- (4.25,-4.15);
  \node[rm note, text=rmBlockLine, anchor=north, align=center] at (4.25,-4.15)
    {privacy wall:\\only fingerprints pass};

  % ---- 6: decide ----------------------------------------------------------------------------
  \node[rm dec, minimum width=27mm] (ovr) at (15.0,-5.45) {override?};
  \draw[rm flow] (jdg) -- (ovr);
  \node[rm stage, minimum width=60mm, minimum height=11mm] (scr) at (8.55,-5.45)
    {\textbf{6\enspace Score}\enspace $L = -6.91 + {}$ strongest weight of each failed question\\
     \tiny e.g.\ pasted \texttt{curl \ldots{} | sh}: $-6.91 + 0.8 + 1.0 + 2.5 = -2.61$};
  \draw[rm flow] (ovr) -- node[rm no, above] {no} (scr);

  % outcomes
  \node[rm allow, minimum width=24mm] (alw) at (2.1,-7.75) {ALLOW};
  \node[rm hold, minimum width=24mm] (hld) at (6.4,-7.75) {HOLD};
  \node[rm block, minimum width=24mm] (blk) at (10.7,-7.75) {BLOCK};
  \draw[rm flow] (scr.south -| blk.north) -- node[rm lbl, align=center] {$L \ge -3.91$ and a\\Contract signal} (blk.north);
  \draw[rm flow] (scr.west) -| node[rm lbl, align=center, pos=0.75] {$L < -4.94$} (alw.north);
  \draw[rm flow] (scr.south -| hld.north) -- node[rm lbl, align=center] {$L \ge -4.94$, or no\\Contract signal} (hld.north);
  \draw[rm flow] (ovr.south) |- node[rm lbl, align=center, pos=0.35] {\textbf{yes}: canary, protected\\asset, attack on Separ\\(trusted source: ALLOW)} (blk.east);

  % ---- 7: gate (the mode) -------------------------------------------------------------------
  \node[rm box, minimum width=88mm, minimum height=10mm, draw=rmDecideLine, fill=rmDecide, thick] (gat) at (8.55,-9.55)
    {\textbf{7\enspace Gate: the mode decides what happens}\\\tiny
     \textbf{observe}: only recorded (``would block'')\quad
     \textbf{warn}: the user is told, nothing stopped\quad
     \textbf{block}: act as below};
  \draw[rm flow] (hld) -- (hld |- gat.north);
  \draw[rm flow] (blk) -- (blk |- gat.north);
  \node[rm term, minimum width=24mm] (run) at (2.1,-9.55) {runs normally\\\tiny (still logged)};
  \draw[rm flow] (alw) -- (run);

  % ---- 8: what happens in block mode --------------------------------------------------------
  \node[rm box, draw=rmHoldLine, fill=rmHold!40, minimum width=34mm, minimum height=12mm] (hep) at (5.3,-11.75)
    {\textbf{hold on a workstation}\\\tiny paused; the user is asked\\\tiny no answer in 60 s = stop};
  \node[rm box, draw=rmHoldLine, fill=rmHold!40, minimum width=34mm, minimum height=12mm] (hsv) at (8.95,-11.75)
    {\textbf{hold on a server}\\\tiny paused and quarantined\\\tiny until the admin decides};
  \node[rm box, draw=rmBlockLine, fill=rmBlock!50, minimum width=40mm, minimum height=12mm] (ctn) at (12.85,-11.75)
    {\textbf{8\enspace Contain} (labelled objects only)\\\tiny kill process tree, cut its network,\\\tiny quarantine new files, undo persistence};
  \draw[rm flow] (gat.south -| hep.north) -- (hep.north);
  \draw[rm flow] (gat.south -| hsv.north) -- (hsv.north);
  \draw[rm flow] (gat.south -| ctn.north) -- (ctn.north);
  \node[rm note, anchor=north] at (10.4,-12.45) {real only with \texttt{separ serve -{}-enforce} as root, otherwise simulated};

  % ---- 9: report and the signed policy ------------------------------------------------------
  \node[rm store, minimum width=25mm, minimum height=19mm] (rep) at (16.3,-12.0)
    {\textbf{9\enspace Report}\\\tiny every decision logged\\\tiny logs: fingerprints only\\\tiny notifications: pseudonym\\\tiny dashboard: \texttt{separ gui}\\\tiny health: \texttt{separ health}};
  \draw[rm ctrl] (ctn.east) -- (rep.west |- ctn.east);
  \node[rm human, minimum width=30mm, minimum height=12mm] (pol) at (15.6,-9.35)
    {\textbf{Signed policy}\\\tiny mode, contracts, weights,\\\tiny trusted sources; signed on\\\tiny the admin workstation};
  \draw[rm back] (pol.west) -- (gat.east |- pol.west);
  \node[rm note, text=rmStageLine, anchor=south] at (pol.north) {read by every stage};
\end{tikzpicture}
}
  \caption{The Separ pipeline end to end (default policy). The thresholds $-4.94$ (hold) and
  $-3.91$ (block) come from the costs in the signed policy; a block also needs a Contract
  signal.}
  \label{fig:flow}
\end{figure}

\subsection{Where it looks}
\begin{tabularx}{\textwidth}{@{}lX@{}}
\toprule
\textbf{Mediation point} & \textbf{When it decides} \\
\midrule
\cmd{separ guard -- <command>} & before the command runs \\
Browser extension & when a page writes to the clipboard, a download starts, optionally a page is opened \\
Watched sandbox and web root & when a file appears \\
Connected sensors (Tetragon, Falco, sshd, Shibboleth, \ldots) & after the event, from their logs \\
Rootkit scan (\cmd{separ scan-rootkit}) & periodically, read-only \\
\bottomrule
\end{tabularx}

\subsection{Honest limits}
\begin{itemize}
  \item \textbf{Pasted install commands are blocked in block mode.} A command pasted from a web
    page that visibly downloads, escalates or pipes into a shell fails all three questions. In the
    test set this blocked \cmd{git clone}, the \cmd{rustup} installer, \cmd{sudo apt update} and
    \cmd{wget}. Listing trusted destinations does not change this.
  \item \textbf{Typed commands that fetch or escalate are held} (asked), never blocked.
  \item \textbf{Not seen:} a terminal paste without the guard or the extension, a script run by name,
    a compiled program, anything not connected.
  \item \textbf{Stopping Separ:} root can stop the system service; a user can stop their own
    per-user daemon. Separ reports its stop and an outside health check notices a missing
    heartbeat, but neither prevents it.
  \item \textbf{Containment} is simulated unless the system service runs as root with
    \cmd{--enforce} (Section~\ref{sec:extras}); even then only processes a sensor reported by
    process ID are touched.
\end{itemize}

\subsection{Separ and Falco}
Falco watches system calls and alerts \emph{after} an action; Separ judges the commands routed to
it \emph{before} they run. They do different jobs and can run together: Separ can read Falco's
alerts as a sensor (Section~\ref{sec:extras}). In a lab test on 22 harmless attack stand-ins,
Falco 0.45.2 with its 25 default rules raised no alert, and with 81 rules (adding its experimental
rule sets) it alerted on 3 of them and on 1 of 6 harmless commands; Separ held or blocked 21 and
none of the harmless ones. \textbf{This test favours Separ}: its author wrote the stand-ins, Falco's
default rules focus on containers, the stand-ins never reach a real network, and everyday commands
were not part of it. It does not show how either tool does against real attacks (details:
\texttt{docs/v2.4-evaluation.md} in the repository).

% ============================================================================
\section{Approvals before installing}
\label{sec:approvals}

\begin{enumerate}
  \item \textbf{IT security officer}: approves the pilot, the machines and the persons with access
    to the logs.
  \item \textbf{Data protection officer}: decides whether a data protection impact assessment under
    Art.~35 GDPR is required, and records the processing in the record of processing activities
    (Art.~30 GDPR). The fact sheet in Section~\ref{sec:facts} lists what is stored.
  \item \textbf{Staff council}: for machines used by staff and for servers staff log in to, because
    Separ records user activity on them.
  \item \textbf{Written notice to users} of the pilot machines before the pilot starts: what is
    recorded, for how long, who can see it, and whom to contact.
  \item \textbf{Private devices}: only with the owner's explicit opt-in.
\end{enumerate}

\subsection{Data protection fact sheet}
\label{sec:facts}
From a lab run of Separ 2.4.0 (\texttt{docs/it-facts.md}, section 4).

\noindent{\small
\begin{tabularx}{\textwidth}{@{}>{\raggedright\arraybackslash}p{0.3\textwidth}X>{\raggedright\arraybackslash}p{0.22\textwidth}@{}}
\toprule
\textbf{File} & \textbf{Data} & \textbf{Form} \\
\midrule
\cmd{logs/events.jsonl} & command, URL, path, file content, login line & keyed fingerprint \\
 & machine of guarded commands, browser, sandbox & pseudonym \\
 & machine of events from connected sensors & \textbf{clear text} host name \\
 & login name (\cmd{actor.user}); account of login sensors & \textbf{clear text} \\
 & source address of a login & class or keyed fingerprint \\
\cmd{logs/verdicts.jsonl} & decision, score, explanation; user names & \textbf{user names clear} \\
\cmd{logs/notifications.jsonl}, webhook & machine, fixed sentences, decision & pseudonym, no user \\
\cmd{logs/barrier.log}, dashboard & quarantined file names; work directory path & \textbf{clear text} \\
\cmd{heartbeat.json} & machine, daemon process ID & pseudonym \\
decoy proposals (only with decoys) & address of an outside / inside source & clear / sealed \\
\bottomrule
\end{tabularx}}

\begin{warnbox}
\textbf{Log permissions.} The system service creates its logs with the default umask in
\cmd{/var/lib/separ} (mode 0755), so in Separ 2.4.0 \textbf{every local user can read them}. Before
the pilot, restrict them with a drop-in: \cmd{sudo systemctl edit separ}, add the two lines
\cmd{[Service]} and \cmd{UMask=0077}, then \cmd{sudo systemctl restart separ} (files created
earlier keep their mode). The logs are not rotated by Separ; delete or archive them at the end of
the pilot as agreed with the data protection officer.
\end{warnbox}

Retention, access and deletion are for the data protection officer to set; Separ itself deletes
nothing. The dashboard shows the last events, so on a server with connected sensors it shows other
users' account names; give dashboard access only to the persons named in the approval.

% ============================================================================
\section{Minimum requirements}
\label{sec:req}
\begin{itemize}
  \item Linux on x86\_64 (\cmd{uname -m} shows \cmd{x86\_64}) with systemd; Ubuntu 24.04, Debian 13,
    or an RHEL-family 9 system with Python 3.12 (installed with the package, see
    Section~\ref{sec:install}). Debian 12 is not supported (Python 3.11).
  \item Root once, to install the package and the signed policy.
  \item No internet access at install time (the package carries its Python dependencies), no
    incoming network: the daemon listens on \cmd{127.0.0.1} only (port 8765 by default).
  \item No kernel module, no cloud account, no licence.
  \item Hardware: \textbf{not measured}. CPU, memory and disk use are to be measured in the pilot.
    A full rootkit scan took about 26 seconds on a development laptop.
\end{itemize}

% ============================================================================
\section{Roles and key custody}
\label{sec:roles}
\begin{description}
  \item[Policy owner] (one URZ administrator and a deputy). Holds the private signing key
    \cmd{admin.key} on an \textbf{admin workstation} (default
    \cmd{\textasciitilde/.separ-admin/admin.key}, mode 0600), \textbf{never on a protected
    machine}. Keeps an offline backup (for example an encrypted USB stick in a safe). Signs every
    policy change.
  \item[Responder.] Reads notifications and the dashboard, answers holds on servers, collects the
    pilot numbers.
\end{description}
The protected machines receive only \cmd{policy.yaml} and \cmd{admin.pub}; the daemon verifies the
signature and cannot sign anything itself.

\paragraph{If the key is lost.} Policies already installed keep working, but no change can be signed.
Create a new key and policy on the admin workstation (Section~\ref{sec:policy}, with a new
\cmd{--key} path and a new \cmd{--config-dir}), install the new \cmd{admin.pub} and \cmd{policy.yaml}
on every machine, and reload.

\paragraph{If the key is stolen.} The thief can sign a policy that switches Separ to observe mode
or allows anything. Separ has no revocation list: replace \cmd{admin.pub} and \cmd{policy.yaml} on
every machine at once (as above), compare the digest that \cmd{separ policy show} prints on each
machine with the one you signed, and report the incident to the IT security officer.

% ============================================================================
\section{Installation}
\label{sec:install}
Download the package and \cmd{SHA256SUMS} of the release from \releases, verify, install:
\begin{lstlisting}
BASE=https://github.com/parsabe/Separ/releases/download/v2.4
curl -fLO $BASE/SHA256SUMS
# Ubuntu 24.04 / Debian 13
curl -fLO $BASE/separ_2.4.0_amd64.deb
sha256sum --ignore-missing -c SHA256SUMS
sudo apt install ./separ_2.4.0_amd64.deb
# RHEL-family 9
curl -fLO $BASE/separ-2.4.0-1.x86_64.rpm
sha256sum --ignore-missing -c SHA256SUMS
sudo dnf install python3.12 ./separ-2.4.0-1.x86_64.rpm
\end{lstlisting}
The release is built and install-tested on Ubuntu 24.04, Debian 13 and Rocky Linux 9 by the
project's CI \todo{confirm that the release v2.4 has been published with its packages}.

\paragraph{Where the files go.} \cmd{/usr/bin/separ}; a private Python environment in
\cmd{/opt/separ/venv} built from the wheels in \cmd{/opt/separ/wheels}; templates in
\cmd{/etc/separ/} (\cmd{ingest.yaml}, examples, \cmd{mappings/}); the system unit
\cmd{/usr/lib/systemd/system/separ.service} and the per-user unit
\cmd{/usr/lib/systemd/user/separ-user.service}, neither enabled; the browser extension and its
policy templates in \cmd{/usr/share/separ/}. \textbf{No policy and no key are installed.} State and
logs go to \cmd{/var/lib/separ}.

\paragraph{Check.}
\begin{lstlisting}
separ --version            # separ 2.4.0
separ health --json        # "not_running" (exit code 1) until the service is started
\end{lstlisting}

% ============================================================================
\section{Policy workflow}
\label{sec:policy}

\subsection{Create the policy on the admin workstation}
\begin{lstlisting}
separ setup --server --mode observe --config-dir ./tubaf-policy \
  --campus-cidr <your campus range> --vpn-cidr <your VPN range> \
  --wlan-cidr <your WLAN range> --mail-domain <your mail domain> \
  --key ~/.separ-admin/admin.key
\end{lstlisting}
Use \cmd{--laptop} instead of \cmd{--server} for workstations. \cmd{--campus-cidr} is required with
\cmd{--server}; Separ has no campus range built in. The command prints exactly what it configured
and writes \cmd{policy.yaml} and \cmd{admin.pub} to \cmd{./tubaf-policy}; the key is created if it
does not exist.

\subsection{Install it on a protected machine}
\begin{lstlisting}
scp tubaf-policy/policy.yaml tubaf-policy/admin.pub <host>:
sudo install -o root -g root -m 0644 policy.yaml /etc/separ/policy.yaml
sudo install -o root -g root -m 0644 admin.pub /etc/separ/admin.pub
separ policy show --policy /etc/separ/policy.yaml     # signature, mode, digest
sudo systemctl enable --now separ.service
\end{lstlisting}
The daemon refuses to start without a signed policy, and refuses a policy or key that is a symlink,
not owned by root, or writable by group or others. Per-user daemons on IT-managed machines use
\cmd{/etc/separ} whenever it holds a policy, under the same rules.

\subsection{Change the mode or the policy}
On the admin workstation, then copy and install as above, then reload:
\begin{lstlisting}
separ policy set-mode warn --policy tubaf-policy/policy.yaml --key ~/.separ-admin/admin.key
separ policy sign --policy tubaf-policy/policy.yaml --key ~/.separ-admin/admin.key  # after editing
sudo separ reload --workdir /var/lib/separ        # or: sudo systemctl reload separ
\end{lstlisting}
Reloading only re-reads and verifies the file; a rejected file leaves the last valid policy in force
and \cmd{separ health} reports \emph{Policy rejected}.

% ============================================================================
\section{Pilot plan}
\label{sec:pilot}

\paragraph{Setup.} 3--5 Linux machines, observe mode, about two weeks, approvals in place, users
informed, log permissions restricted (Section~\ref{sec:facts}). Users who want their pasted commands
judged run them through \cmd{separ guard -- <command>}; with the browser extension (optional) page
clipboard writes are judged without that.

\paragraph{Each week, per machine.}
\begin{lstlisting}
sudo grep -c 'Would hold (not enforced)'  /var/lib/separ/logs/notifications.jsonl
sudo grep -c 'Would block (not enforced)' /var/lib/separ/logs/notifications.jsonl
\end{lstlisting}
Record: the number of would-hold and would-block stories; for each, what it was (the notification and
the dashboard explain it in sentences; Separ does not store the command text, so ask the user what
they ran at that time); and whether it was a false alarm. Note also the CPU and memory use of the
\cmd{separ} process and the size of \cmd{/var/lib/separ/logs}.

\paragraph{The dashboard.} \cmd{sudo separ gui --workdir /var/lib/separ --no-browser} prints the
local address with its token; open it on the machine itself (or through an SSH tunnel to port
8765). It shows each story with its three questions, score and explanation.

\paragraph{Criteria to move on} (to be fixed with the IT security officer before the pilot):
\begin{itemize}
  \item to \textbf{warn} mode (users see a notice, nothing is stopped): no would-block on routine
    work that cannot be fixed by policy, and at most about one would-hold per user and week;
  \item to \textbf{block} mode: a further period in warn mode meeting the same targets, a decision
    on pasted install commands (they will be blocked), and an agreed response process.
\end{itemize}

% ============================================================================
\section{Optional extras}
\label{sec:extras}
Each extra has its own requirement and is not needed for the pilot.
\begin{description}
  \item[Browser extension.] Requires Firefox ESR 128+ or Chromium/Chrome 121+ and a per-user daemon.
    As root (for Chromium add \cmd{--chromium-extension-id <id>}):
\begin{lstlisting}
sudo separ browser-install --system
\end{lstlisting}
    then roll out the extension with the enterprise policy templates in
    \path{/usr/share/separ/browser/}.
  \item[Real enforcement.] Requires the system service as root, cgroup v2 and \cmd{nft} for the
    network cut, and a sensor that reports process IDs (Tetragon or Falco). Add \cmd{--enforce} to
    the \cmd{ExecStart} line of the unit with \cmd{sudo systemctl edit separ}. \cmd{separ health}
    then lists each containment action as available or not.
  \item[Kernel sensors] (Tetragon, Falco). Install the sensor, enable its mapping in
    \path{/etc/separ/ingest.yaml}, and test a sample first:
\begin{lstlisting}
separ ingest list
separ ingest check /etc/separ/mappings/<mapping>.yaml --input <sample> --judge
\end{lstlisting}
  \item[Login and mail sensors] (sshd, Shibboleth IdP, Keycloak, Postfix/Dovecot, M365 audit). Their
    rules start in observe mode and a story only they produced is held at most. They record account
    names in clear text; include them in the data protection decision.
  \item[Notifications.] Options of \cmd{separ serve} (add them to \cmd{ExecStart}); \cmd{--notify-on}
    chooses the categories. Messages name the machine by pseudonym only.
\begin{lstlisting}
--notify-webhook <url>          # a Mattermost or Slack incoming webhook
--notify-command "mail -s Separ admin@example.org"
\end{lstlisting}
\end{description}

% ============================================================================
\section{Operations}
\label{sec:ops}

\paragraph{Health check and watchdog.} \cmd{separ health} reads the heartbeat file: exit code 0
healthy, 1 not running or stale, 2 running with a problem. A cron entry (\cmd{crontab -e} as root)
notifies once per outage and once on recovery:
\begin{lstlisting}
*/5 * * * * separ health --workdir /var/lib/separ --notify-command "mail -s Separ admin@example.org"
\end{lstlisting}

\paragraph{Rootkit scans.} The service scans every hour (\cmd{--rootkit-interval}, default 3600 s),
sweeping only the process IDs allocated so far. A full manual scan:
\cmd{sudo separ scan-rootkit} (about 26 s on a development laptop; \cmd{--pid-sweep bounded} is
faster). Measure the transient rate once on a clean machine:
\cmd{sudo separ scan-rootkit --calibrate 20 --workdir /var/lib/separ}. Exit codes: 0 nothing to
review, 1 findings to review, 2 block-level findings.

\paragraph{Updating.} Install the new release's package as in Section~\ref{sec:install}, then
\cmd{sudo systemctl restart separ}. The policy in \cmd{/etc/separ} and the state in
\cmd{/var/lib/separ} are kept.

\paragraph{Uninstalling and rolling back.} \cmd{sudo apt remove separ} or \cmd{sudo dnf remove
separ} removes the program, the private environment and the files the package installed in
\cmd{/etc/separ}; the policy, the public key and \cmd{/var/lib/separ} (logs) stay and must be
deleted or archived as agreed. To roll back, remove the package and install the previous release's
package file.

\paragraph{Troubleshooting.}
\begin{description}
  \item[``no signed policy at \ldots: run \cmd{separ setup} first''] Install \cmd{policy.yaml} and
    \cmd{admin.pub} in \cmd{/etc/separ} (Section~\ref{sec:policy}).
  \item[``refusing \ldots: writable by group or others'' / ``not root''] Reinstall the files with
    \cmd{install -o root -g root -m 0644}; the folder \cmd{/etc/separ} must also be root-owned and
    not group- or world-writable.
  \item[Policy rejected (\cmd{unsigned}, \cmd{bad\_signature})] The file was edited without signing,
    or signed with another key. Sign it on the admin workstation and install it again.
  \item[Degraded protection] Lists what Separ cannot see on this machine (for example no kernel
    sensor, simulated containment). Normal for the pilot; it is information, not an error.
\end{description}

% ============================================================================
\section{FAQ for users and service desk staff}
\label{sec:faq}
\begin{description}
  \item[Why was my command held?] It came from outside (for example pasted from a web page) or did
    something its program normally does not, and it looked like it could download, hide, persist or
    reach credentials. In observe mode nothing was stopped; in block mode a held command run through
    \cmd{separ guard} did not run.
  \item[Why was a normal install command blocked?] Commands pasted from a web page that download
    and run code look exactly like the ``fix your browser'' lures. Typed in by hand they are only
    held. During the pilot nothing is blocked.
  \item[How do I report a false alarm?] Tell the service desk the time, the machine and what you
    were doing; the responder finds the story in the dashboard. A policy owner can add a narrow,
    expiring exception to the policy and sign it.
  \item[Does Separ read my commands or files?] It keeps no command text, addresses or file contents,
    only keyed fingerprints. It does record your login name and the time of each event.
  \item[Who do I contact?] The URZ service desk; for the project: \contact.
\end{description}

\end{document}
